\documentclass[11pt]{amsart}

\usepackage[T1]{fontenc}
\usepackage[utf8]{inputenc}
\usepackage{lmodern}
\input{glyphtounicode}
\pdfgentounicode=1

\usepackage{amsmath,amssymb}
\usepackage{array,booktabs,tabularx}
\usepackage{geometry}
\usepackage{xcolor}
\usepackage{enumitem}
\usepackage{url}
\usepackage{hyperref}
\geometry{margin=1in}
\setlength{\emergencystretch}{3em}
\hypersetup{
  colorlinks=true,
  linkcolor=blue!45!black,
  citecolor=blue!45!black,
  urlcolor=blue!45!black,
  pdftitle={Software and Reproducibility for Numerical Experiments on Projected Multilinear Models},
  pdfauthor={Eliuth Chavero Jasso},
  pdfsubject={Reproducibility, provenance and numerical controls for two independent computational studies},
  pdfkeywords={reproducible research, numerical software, provenance, interval arithmetic, CPU/GPU agreement, experiment identity}
}

\newcommand{\norm}[1]{\lVert#1\rVert}
% \path is provided by the url package loaded above; do not redefine it.

\theoremstyle{remark}
\newtheorem{remark}{Remark}[section]

\title[Software and reproducibility for projected multilinear models]{Software and Reproducibility for Numerical Experiments\\ on Projected Multilinear Models}
\author{Eliuth Chavero Jasso}
\address{Independent Researcher, Apizaco, Tlaxcala, Mexico}
\thanks{This article makes no independent mathematical claim. Every mathematical statement
it refers to belongs to \cite{ChaveroJassoTrees} or \cite{ChaveroJassoNumerical}, and is
cited there.}
\date{31 July 2026}

\begin{document}
\maketitle

\begin{abstract}
We describe the software supporting two independent computational studies of multilinear
maps under orthogonal projection, and the measures taken to make their reported results
reconstructible. The first study evaluates finite composition trees of multilinear maps
whose intermediate outputs are recursively projected, and computes upper bounds and
rigorous interval enclosures for the resulting error \cite{ChaveroJassoTrees}. The second
fits cyclically symmetric multilinear maps together with orthogonal projectors and measures
fourteen diagnostic quantities \cite{ChaveroJassoNumerical}. The two share no code and no
evidence; they are presented here as two case studies of one set of reproducibility
practices, not as one system.

We describe: how multilinear maps, typed trees and projectors are represented; why each
study carries a reference implementation alongside an optimised one, and what is compared
between them; how numerical precision is controlled, including a process-global
reduced-precision setting that silently defeats such control and how it is contained; how a
mathematical configuration is distinguished from an execution of it, so that repeated and
resumed executions are never counted as independent observations; what each run records;
and how the tables, figures and manuscripts are regenerated. We report what has been
verified --- $73$ and $85$ automated tests, agreement between reference and optimised
implementations, agreement between two devices in double precision --- and what has not:
in particular, no reproduction of either study from a clean environment by an independent
party has been performed.
\end{abstract}

\subjclass[2020]{68N30, 65Y05, 65G50, 15A69}
\keywords{reproducible research, numerical software, provenance, interval arithmetic,
CPU/GPU agreement, experiment identity}

\section{Scientific purpose of the software}
\label{sec:purpose}

Two questions are being asked, by two disjoint bodies of code.

\paragraph{Case study I: bounds for recursively projected composition trees.}
Given a finite tree whose internal vertices carry multilinear maps between
finite-dimensional Hilbert spaces, and orthogonal projectors onto prescribed subspaces at
every vertex, how large can the discrepancy between unprojected and recursively projected
evaluation be? The software enumerates trees exactly, evaluates both quantities, computes
four families of upper bound, and produces rigorous interval enclosures of explicit lower
constructions. Its scientific role is to \emph{check} the estimates of
\cite{ChaveroJassoTrees} on many configurations, and to \emph{search} for constructions that
approach them. It cannot prove anything, and no bound it computes is a proof.

\paragraph{Case study II: diagnostics for fitted multilinear maps and projectors.}
Given a cyclically symmetric multilinear map in a low-rank representation together with a
fitted orthogonal projector, what do fourteen diagnostic quantities --- projector
identities, a commutator residual, closure and associator defects, principal angles between
subspaces at different resolutions, mode ranks, and a generalised Jacobi residual --- take
as values? Its scientific role is to \emph{measure}, and in four cases the measurement is a
negative result \cite{ChaveroJassoNumerical}.

\begin{remark}[The two case studies are kept apart]
\label{rem:separate}
No module is shared, no dataset is shared, and no result from one is used as evidence for
the other. They appear in one article because they are governed by one set of
reproducibility practices, and because a reader wishing to rebuild either will follow the
same steps. Any future connection between them would require a theorem, and none exists.
\end{remark}

\section{Mathematical objects and their representation}
\label{sec:objects}

\subsection{Case study I}

\begin{center}
\small
\begin{tabularx}{\linewidth}{@{}lX@{}}
\toprule
object & representation \\
\midrule
typed space & ambient dimension, reduced dimension, isometry $Q$, projector $P=QQ^{*}$; the identities $P^{2}=P$, $P^{*}=P$ and $Q^{*}Q=I$ are validated on construction \\
multilinear map & ordered input signature, output type, and a dense coefficient tensor whose shape must agree with the signature \\
leaf & a label and a type; the label must resolve to reduced input data of the declared size \\
internal vertex & a map identifier, an output type, and ordered children whose output types must equal the corresponding input types \\
tree & an immutable structure with a canonical serialisation and a content hash \\
\bottomrule
\end{tabularx}
\end{center}

Validation precedes evaluation. An edge whose types do not match is rejected before any
arithmetic occurs; it is a control, not a numerical failure. Forty-eight such
configurations are declared and all are rejected.

\subsection{Case study II}

A multilinear map is stored in a rank-$R$ canonical polyadic representation
\cite{KoldaBader2009}, which avoids materialising a coefficient tensor of size
$n^{a+1}$. The projector is $P=UU^{*}$ for an isometry $U\in\mathbb C^{n\times r}$. Cyclic
symmetrisation is applied by an orthogonal projector onto the cyclically invariant subspace.
A reduced tensor is obtained by contracting the map against $U$ in every slot.

\begin{remark}[Representation is not structure]
\label{rem:representation}
A low-rank representation is a storage device. It is not unique --- factors may be rescaled
subject to a product constraint and permuted --- so any comparison of factors must be
performed modulo that freedom. This is why the comparison tooling of
\cite{ChaveroJassoNumerical} reports several aligned quantities separately rather than one
score.
\end{remark}

\section{Reference and optimised implementations}
\label{sec:implementations}

Both case studies carry two implementations of their core computation, and compare them.
The purpose is to detect errors that a single implementation cannot reveal.

\paragraph{Case study I.} A reference evaluator applies each coefficient tensor with
explicit index loops, in the most direct transcription of the definition. An optimised
evaluator contracts modes through vectorised operations \cite{HarrisEtAl2020}, and a third
runs on a GPU \cite{PaszkeEtAl2019}. All three are exercised on real and complex data, on
arities two through four, with repeated leaves, and against exactly invariant controls for
which the correct answer is known in closed form.

\paragraph{Case study II.} A reference implementation of the reduced tensor uses explicit
loops; an optimised one uses a contraction exploiting the low-rank structure, including the
cyclic average. They agree to better than $10^{-10}$ in double precision and $10^{-4}$ in
single precision, and an exact rational case at $n=2$, $r=1$ is checked against a hand
computation with no floating-point arithmetic involved.

\paragraph{Fresh implementation rather than reuse.} The core numerical primitives of case
study II were reimplemented rather than imported from the historical code, for the reason
given in Section~\ref{sec:precision}. The reimplementation is checked against the
historical formulas by direct citation of the source lines it reproduces.

\section{Numerical precision and hardware controls}
\label{sec:precision}

\subsection{A process-global setting that defeats precision control}

The historical code sets, unconditionally and at import time when a GPU is present, two
process-global flags: one enabling reduced-precision matrix multiplication on the GPU
(TF32), and one lowering the default precision of single-precision matrix multiplication.
Neither is scoped to the importing module.

The consequence is that a process which imports that module for any reason has its own
precision discipline silently overridden, whatever it set beforehand. This is not a
hypothetical: it is why the primitives of case study II were reimplemented rather than
imported. The reimplementation exposes a single idempotent entry point that establishes the
required precision settings, together with an assertion that re-checks them, so that the
discipline can be re-established at any point rather than assumed.

\subsection{A second hazard: moving a model between devices}

Transferring a model to another device with the framework's standard method updates the
tensors but not plain attributes recording the intended device and data type. Code reading
those attributes then disagrees with the tensors. Device transfers are therefore performed
by constructing an instance on the target device and copying parameter values explicitly,
not by the standard method.

\subsection{What is compared, and how}

Both case studies compare a CPU and a GPU execution of the same configuration, in double
precision, with reduced-precision matrix multiplication explicitly disabled and
deterministic algorithms enabled.

\begin{center}
\small
\begin{tabularx}{\linewidth}{@{}lrrX@{}}
\toprule
quantity & CPU & GPU & note \\
\midrule
idempotence residual & $1.312\times10^{-15}$ & $1.198\times10^{-15}$ & both at the noise floor \\
self-adjointness residual & $0$ (exact) & $2.220\times10^{-16}$ & both at the noise floor \\
unexplained relative residual & $1.0000380759060201$ & $1.0000380759060197$ & relative difference $4.4\times10^{-16}$ \\
\bottomrule
\end{tabularx}
\end{center}

Only the third quantity is large enough for a relative comparison to mean anything; the
first two are at the noise floor on both devices, so their agreement establishes that
neither device is producing something structurally wrong, and no more. For case study I the
largest observed difference between the CPU and GPU evaluators over the whole registry was
$1.922\times10^{-8}$, which is a validation threshold on the implementation and not a
quantity appearing in any theorem.

\begin{remark}[Timing is not compared here]
\label{rem:timing}
A single timing comparison at one problem size is not a hardware-performance result. The
throughput question is answered instead by the swept measurement of
\cite{ChaveroJassoNumerical}, over $160$ configurations and four ambient dimensions, and
even that is a measurement on one machine.
\end{remark}

\section{Experiment specification}
\label{sec:specification}

An experiment is specified by a declarative configuration: the arity, the ambient
dimension, the reduced rank, the representation rank, the seed, the device, the numerical
precision, the objective, the number of optimisation steps, and the tolerances. The
configuration is resolved and validated before execution, and the resolved form --- not the
form as written --- is what is recorded.

\subsection{Configuration, execution, restart}

This distinction is the one that most affects how results may be reported.

\begin{center}
\small
\begin{tabularx}{\linewidth}{@{}lX@{}}
\toprule
identity & meaning \\
\midrule
configuration & one distinct mathematical setting. Two runs with the same arity, dimension, ranks and seed are the same configuration, whatever hardware they used \\
execution & one run of one configuration. Distinguished by device, by time, and by whether it resumed from a checkpoint \\
optimiser restart & one search trajectory within an execution \\
\bottomrule
\end{tabularx}
\end{center}

The ledger records both identifiers for every run, together with the parent execution when
a run resumed from a checkpoint. This is what makes the following statements checkable
rather than assertions: that the $416$ executions of the swept extension in
\cite{ChaveroJassoNumerical} cover $208$ configurations, each executed once per device; and
that the nineteen historical runs of the same study form one chain of resumptions rather
than nineteen independent trials.

\begin{remark}[Why this matters for reporting]
\label{rem:accounting}
Counting executions as configurations inflates an apparent sample size by whatever factor
the execution grid happens to have. Counting optimiser restarts as replicates does the same
and additionally suggests an uncertainty estimate that does not exist. Both are prevented
here by the identity scheme rather than by discipline in writing.
\end{remark}

\subsection{Work that was specified but not executed}

Case study I specifies an extended optimiser grid of $460{,}800$ search trajectories that
was not executed, under a stated computational-cost constraint. Case study II specifies two
further sweep stages that were not executed. In both, the unexecuted work is recorded in
the configuration and marked as not executed. It is not omitted, and the executed portion is
never presented as though it were the whole.

\section{Provenance}
\label{sec:provenance}

\subsection{What a run records}

Every run directory contains: the resolved configuration; the exact command; the software
environment and dependency versions; the hardware; the source commit; content hashes of the
mathematical objects and of the input data; the input tensors; the computed metrics; any
enclosures; the logs; and hashes of every output file. History is append-only: a retry
increments a counter and adds a record, and never rewrites an earlier one.

\subsection{Historical evidence is ingested, not trusted}

The historical evidence for case study II was ingested by a tool that hashes each source
file, hashes each run directory, parses the run log, cross-references it against the
on-disk summaries, reconstructs the chain of resumptions, and reclassifies every historical
run against the current vocabulary. That process is what established the provenance
findings reported in \cite{ChaveroJassoNumerical}: that every historical run was
exploratory despite directory names suggesting otherwise, that nine logged runs form a
single chain, and that the analysis script changed mid-campaign without this being
recorded.

\subsection{A corrected classification, re-derived rather than re-run}

The classification of enclosure outcomes in case study I was found to be defective: the
label reserved for an exactly determined constant was assigned only when the proved upper
bound was itself zero --- the vacuous case --- while configurations whose lower and upper
bounds genuinely coincided at a positive value received a weaker label, and configurations
for which no positive lower bound had been obtained were described as having a certified
lower bound.

The classification is a pure function of the certified lower and upper bounds, both of
which are recorded. It was therefore corrected in the source and \emph{re-derived} from the
recorded bounds, without re-running any experiment. The re-derivation is a separate script,
its output is written alongside the original rather than over it, and the correspondence
between the old and new labels is a bijection on the agreement classes, so the two
derivations can be compared directly. Historical outputs were not modified.

\section{Reconstructing the tables, figures and manuscripts}
\label{sec:reconstruct}

\subsection{Tables}

Tables in \cite{ChaveroJassoTrees} are generated from the recorded registries by a script
that also emits the numerical macros used in the text, so that no number in the manuscript
is typed by hand. Where the generated text carried implementation vocabulary, a separate
and auditable rewriting step replaces it with standard terminology; that step touches
headings and labels only, never a numeric value, and it refuses to complete if any
forbidden token survives.

\subsection{Figures}

Figures in \cite{ChaveroJassoSupplement} are produced by a single generator from committed
data files. It emits each figure as a vector PDF, an SVG and a 300-dpi PNG, and writes a
manifest recording, per figure: the input files with checksums, the output files with
checksums, the checksum of the generator itself, and the versions of the numerical
libraries. No figure value is entered by hand.

\subsection{Manuscripts}

Each manuscript builds from a clean directory with a standard \LaTeX{} toolchain. The build
is checked for undefined references, undefined citations, overfull lines and errors, and the
resulting PDF is checked for correct text extraction --- specifically, that ligatures are
recoverable, which requires an outline font rather than the bitmap fonts that a bare
\path{T1} font encoding selects by default.

\section{Testing strategy}
\label{sec:testing}

\begin{center}
\small
\begin{tabularx}{\linewidth}{@{}lrX@{}}
\toprule
suite & tests & scope \\
\midrule
case study I & $73$ & typed validation, enumeration, evaluator agreement, bound soundness, ordering rule against all permutations through arity seven, interval enclosures, run-artifact contract \\
case study II & $85$ & projector identities, block-level diagnostics, gauge invariance, ledger semantics including retry and checkpoint lineage, and the enforcement point below \\
\bottomrule
\end{tabularx}
\end{center}

Both suites pass. Two categories of test deserve mention.

\paragraph{Negative controls that must fail.} Invalid typed edges must be rejected; a
deliberately sign-mutated identity must be detected; an objective claimed to be invariant
under a change of basis must be checked against an actual change of basis. Each of these
found a real defect during development, which is the reason they exist.

\paragraph{Enforcement rather than convention.} One function is the single point at which a
run's evidential status is assigned, and it raises an exception --- rather than silently
downgrading --- if an exploratory run is ever assigned a status reserved for a rigorous one.
A schema-drift detector pinned to file checksums was verified by deliberately editing a
schema, confirming the test failed, restoring the file, and confirming it passed again.

\begin{remark}[What passing tests establish]
\label{rem:testscope}
Passing tests establish consistency between the implementation and its specification. They
do not establish any mathematical statement in either companion article, and no bound
computed by this software is a proof of anything.
\end{remark}

\section{Cross-platform and cross-implementation agreement}
\label{sec:agreement}

Three kinds of agreement are checked and are reported separately:

\begin{enumerate}[leftmargin=*]
\item between a reference implementation and an optimised one, on the same device
(Section~\ref{sec:implementations});
\item between two devices, on the same configuration in double precision
(Section~\ref{sec:precision});
\item between two independent implementations of the same mathematical quantity written
against different formulations (case study II's reduced tensor).
\end{enumerate}

For each, we report the quantity compared, the precision, the tolerance, the absolute
difference and, where the magnitudes make it meaningful, the relative difference. Where a
quantity is at the numerical noise floor on both sides, we say so and draw no further
conclusion.

\section{Limitations}
\label{sec:limitations}

\begin{enumerate}[leftmargin=*]
\item \textbf{No clean-environment reproduction has been performed.} Neither study has been
rebuilt and re-run from a fresh environment as a dedicated exercise, by the author or by
anyone else. Continuous integration does exercise a clean install of the packaged source,
but it does not then run either study's full experiment set inside that environment. This
is the largest outstanding gap.
\item \textbf{No independent party has reproduced anything.} All the agreement reported
above is internal.
\item Numerical results come from one machine and one software stack. Timing results in
particular do not transfer.
\item The extended optimiser grid of case study I and the two further sweep stages of case
study II were specified and not executed.
\item Multilinear operator norms are computed exactly only for declared explicit
constructions; elsewhere a Frobenius upper bound is substituted, which inflates the
constants that depend on it.
\item Ten of the fourteen experiments of case study II were not extended over a parameter
grid and were not run on a GPU at all.
\item The numerical evidence of case study I was produced at an earlier source commit than
the manuscript text. The provenance file records both, and the corrected classification of
Section~\ref{sec:provenance} was re-derived from the recorded bounds rather than
regenerated by re-running the pipeline at the later commit.
\end{enumerate}

\section{Instructions for independent reproduction}
\label{sec:instructions}

The steps below are what a third party would follow. They have been exercised individually
but not as one continuous pass from a fresh environment; see Section~\ref{sec:limitations}.

\begin{enumerate}[leftmargin=*]
\item \textbf{Environment.} Create an isolated Python environment at the version recorded in
the provenance file and install the locked dependency set. A GPU is optional; every result
except the device-agreement measurements can be reproduced without one.
\item \textbf{Tests.} Run both test suites and confirm the counts in
Section~\ref{sec:testing}. A failure here means the environment differs from the recorded
one, and nothing further should be attempted until it is resolved.
\item \textbf{Data.} Either regenerate the experiment records by running the pipelines with
the recorded configurations, or use the records shipped with the package. Regeneration is
the stronger check and takes substantially longer; the shipped records carry checksums so
that the two can be compared.
\item \textbf{Tables and macros.} Run the table generator and the classification
re-derivation. Both are deterministic and idempotent: running either twice produces
byte-identical output.
\item \textbf{Figures.} Run the figure generator. It writes a manifest with checksums, which
should be compared against the shipped manifest.
\item \textbf{Manuscripts.} Build each manuscript from a clean directory and confirm zero
undefined references, zero undefined citations, zero overfull lines and zero errors, and
confirm that text extraction from the resulting PDF recovers ligatures correctly.
\item \textbf{Checksums.} Verify the package checksum file.
\end{enumerate}

\section{Statement on the use of AI tools}
\label{sec:ai}

An AI assistant was used in preparing the software and the manuscripts: for drafting and
editing prose, for restructuring the manuscripts, for locating defects in the classification
code and in reported numerical values, and for writing auxiliary scripts. All mathematical
statements, all proofs, all citations, all interpretations and the final wording are the
author's responsibility, and the author has checked them. No AI system is an author. Any
citation suggested by such a tool was verified against the primary source before being
retained. Authors submitting elsewhere should consult the target venue's current policy on
AI-assisted writing, which this statement is intended to satisfy but may not, depending on
the venue.

\bibliographystyle{amsplain}
\bibliography{references}

\end{document}
